\section*{Capítulo \ref{ch:b1:e-ph-diagrams} --- Diagramas E--pH}

\begin{solution}{exo:b1:e-ph-diagrams:1}
Manganeso: \ce{Mn} (0) $<$ \ce{Mn^2+}, \ce{Mn(OH)2} ($+$II) $<$ \ce{MnO2}
($+$IV) $<$ \ce{MnO4-} ($+$VII). Cloro: \ce{Cl-} ($-$I) $<$ \ce{Cl2} (0)
$<$ \ce{HClO}, \ce{ClO-} ($+$I) $<$ \ce{ClO3-} ($+$V).
\end{solution}

\begin{solution}{exo:b1:e-ph-diagrams:2}
$-0.059 \times 3/1 = -0.177$; $+0.059 \times 2/2 = +0.059$;
$-0.059 \times 8/5 = -0.094$; $-0.059 \times 2/2 = -0.059$ V por unidad de
pH. En el par del cobre, los protones se liberan del lado de la forma
reducida ($m = -2$): el potencial sube con el pH.
\end{solution}

\begin{solution}{exo:b1:e-ph-diagrams:3}
\ce{H+}/\ce{H2} pasa por 0~V a pH 0 y nunca alcanza 0.50~V para
$\mathrm{pH} \geqslant 0$. \ce{O2}/\ce{H2O} alcanza 0.50~V en
$\mathrm{pH} = 0.73/0.059 = 12.4$, y 0~V solo a pH 20.8, fuera de la
escala. La banda tiene \qty{1.23}{V} de anchura a cualquier pH, incluido
el pH 7 (de $-0.41$ a $+0.82$~V).
\end{solution}

\begin{solution}{exo:b1:e-ph-diagrams:4}
(1, 1.0~V): \ce{Fe^3+}. (4, 0): la recta \ce{Fe(OH)3}/\ce{Fe^2+} está en
$1.09 - 0.71 = 0.38$~V y el punto queda por debajo de ella y por encima
de $-0.47$~V:
\ce{Fe^2+}. (10, $-0.4$~V): entre $-0.07 - 0.59 = -0.66$ y $0.28 -
0.59 = -0.31$~V: \ce{Fe(OH)2}. (10, 0.5~V): \ce{Fe(OH)3}.
\end{solution}

\begin{solution}{exo:b1:e-ph-diagrams:5}
\ce{Fe^2+}/\ce{Fe}: $-0.41 + 0.030 \times (-6) = -0.59$~V.
\ce{Fe^3+}/\ce{Fe(OH)3}: $14.00 - \frac13(38.55 - 6) = 3.15$. Crecen los
dominios de los iones disueltos (corrosión): \ce{Fe^2+} se extiende hasta
$-0.59$~V, y \ce{Fe^3+}, hasta pH 3.15; el metal y los hidróxidos pierden
terreno.
\end{solution}

\begin{solution}{exo:b1:e-ph-diagrams:6}
\ce{Fe(OH)3 + 3H+ + e- -> Fe^2+ + 3H2O}: $E = E^\circ + 0.059\log\frac{[\ce{H+}]^3}{[\ce{Fe^2+}]}
= (E^\circ - 0.059\log c) - 0.177\,\mathrm{pH}$. En pH 1.82, la recta
encuentra \ce{Fe^3+}/\ce{Fe^2+} a 0.77~V, de modo que la constante es
$0.77 + 0.177 \times 1.82
= 1.09$~V.
\end{solution}

\begin{solution}{exo:b1:e-ph-diagrams:7}
$E^\circ(\ce{Cu+}/\ce{Cu}) = 0.52 > E^\circ(\ce{Cu^2+}/\ce{Cu+}) = 0.16$:
según la
\cref{prop:b1:e-ph-diagrams:disproportionation-criterion} \ce{Cu+}
se dismuta. Frontera única: $E = 0.34 + 0.030\log 0.010 =
0.28$~V.
\end{solution}

\begin{solution}{exo:b1:e-ph-diagrams:8}
A pH 0, el dominio del cobre (por debajo de 0.28~V) se solapa con el del
agua y está por encima de la recta del hidrógeno: no hay reacción con el
ácido clorhídrico sin aire. Ácido nítrico: \ce{NO3-}/\ce{NO}, a 0.96~V,
está por encima del dominio del cobre, que se oxida:
\ce{3Cu + 2NO3- + 8H+ -> 3Cu^2+ + 2NO + 4H2O}.
% ledger: dfg:NO3-_ao, dfg:NO_g (E0 computed in figdata/bachelor-1/standard-potentials.py)
\end{solution}

\begin{solution}{exo:b1:e-ph-diagrams:9}
A pH 3, la recta \ce{O2}/\ce{H2O} está en $1.23 - 0.18 = 1.05$~V, por
encima de todo el dominio de \ce{Fe^2+}: el oxígeno oxida el hierro(II).
Por encima de pH 1.82, el hierro(III) formado precipita como hidróxido
pardo amarillento:
\ce{4Fe^2+ + O2 + 10H2O -> 4Fe(OH)3 + 8H+}.
\end{solution}

\begin{solution}{exo:b1:e-ph-diagrams:10}
Verticales: $14.00 - \frac12(16.38 - 2) = 6.81$; y, a partir de
\ce{Zn(OH)2 + 2OH- <=> Zn(OH)4^2-} ($K = 10^{-1.93}$) con
$[\ce{Zn(OH)4^2-}] = 10^{-2}$, $[\ce{OH-}]^2 = 10^{-0.07}$, pH 13.97.
\ce{Zn^2+}/\ce{Zn}: $-0.76 + 0.030 \times (-2) = -0.82$~V. \ce{Zn(OH)2}/\ce{Zn}:
pendiente $-0.059$, que pasa por $(6.81, -0.82)$:
$E = -0.42 - 0.059\,\mathrm{pH}$.
\ce{Zn(OH)4^2- + 4H+ + 2e- -> Zn + 4H2O}: pendiente $-0.118$, que pasa
por $(13.97, -1.24)$: $E = 0.41 - 0.118\,\mathrm{pH}$.
\end{solution}

\begin{solution}{exo:b1:e-ph-diagrams:11}
El dominio del hierro está por debajo de la recta del hidrógeno a
cualquier pH. A pH 2, el hierro se oxida a \ce{Fe^2+} con desprendimiento
de hidrógeno (corrosión). A pH 7, la reacción sigue siendo posible, pero
la recta del hidrógeno ($-0.41$~V) está cerca de \ce{Fe^2+}/\ce{Fe} y la
reacción es muy lenta. A pH 13, el hierro se oxida a \ce{Fe(OH)2}, un
\omterm{def:b1:e-ph-diagrams:corrosion-domains}{dominio de pasivación}: se forma una capa y el clavo conserva su brillo.
Con oxígeno disuelto, la recta \ce{O2}/\ce{H2O}, muy por encima, oxida el
hierro a hierro(III) a cualquier pH: la herrumbre, \ce{Fe(OH)3} y sus
óxidos deshidratados.
\end{solution}

\begin{solution}{exo:b1:e-ph-diagrams:12}
El cinc tiene el dominio más bajo: en contacto con el acero, se oxida
primero, \ce{2Zn + O2 + 2H2O -> 2Zn(OH)2} a pH 8, y los electrones que
libera reducen el oxígeno en la superficie del acero, que permanece
metálico. El magnesio ($E^\circ = -2.36$~V) está aún más abajo y también
sirve; el cobre (0.34~V) está por encima del hierro: unido al cobre, el
hierro sería el metal oxidado, y más deprisa que solo.
\end{solution}

\begin{solution}{pb:b1:e-ph-diagrams:1}
\textbf{1.} $-$I, 0, $+$I, $+$I.
\textbf{2.} \ce{Cl-} abajo, \ce{Cl2}, y después \ce{HClO} y \ce{ClO-}
arriba.
\textbf{3.} \ce{HClO} y \ce{ClO-}, un \omterm{def:b1:predominant-reaction:bronsted}{par ácido--base}.
\textbf{4.} \ce{Cl2 + 2e- -> 2Cl-}.
\textbf{5.} \ce{2HClO + 2H+ + 2e- -> Cl2 + 2H2O}.
\textbf{6.} \ce{ClO- + H2O + 2e- -> Cl- + 2OH-}.
\textbf{7.} \ce{HClO} por debajo de pH 7.55, \ce{ClO-} por encima.
\textbf{8.} No se intercambia ningún electrón; la frontera la fija solo
$K_a$.
\textbf{9.} $1/(1 + 10^{7.4 - 7.55}) = 0.59$: \qty{59}{\percent}.
\textbf{10.} $1/(1 + 10^{0.45}) = 0.26$: tres cuartas partes del cloro
activo serían \ce{ClO-}, el desinfectante más débil.
\textbf{11.} \ce{ClO-}.
\textbf{12.} $E = 1.396 + 0.0296\log\frac{[\ce{Cl2}]}{[\ce{Cl-}]^2} = 1.396
+ 0.0296\log\frac{0.005}{10^{-4}} = 1.396 + 0.050 = 1.446$~V; no hay
\ce{H+} en la \omterm{def:b1:nernst:couple}{semiecuación}.
\textbf{13.} $E = 1.594 + 0.0296\log\frac{[\ce{HClO}]^2[\ce{H+}]^2}{[\ce{Cl2}]}
= 1.594 + 0.0296\log\frac{10^{-4}}{0.005} - 0.0592\,\mathrm{pH} = 1.594 -
0.050 - 0.0592\,\mathrm{pH}$.
\textbf{14.} Igualando, $0.0592\,\mathrm{pH} = 1.594 - 1.396 - 2 \times 0.0503
= 0.0974$, $\mathrm{pH} = 1.646$.
\textbf{15.} \ce{HClO + H+ + 2e- -> Cl- + H2O}: un protón por dos
electrones, pendiente $-0.0592/2 = -0.0296$.
\textbf{16.} Sumando las dos \omterm{def:b1:nernst:couple}{semiecuaciones} (un electrón por átomo de
cloro en cada una): $2E^\circ = 1.594 + 1.396$, $E^\circ(\ce{HClO}/\ce{Cl-}) =
1.495$~V; frontera $E = 1.495 - 0.0296\,\mathrm{pH}$.
\textbf{17.} \ce{ClO- + 2H+ + 2e- -> Cl- + H2O}: $E = 1.718 - 0.0592\,\mathrm{pH}$
(constante elegida por continuidad: $1.495 - 0.0296 \times 7.55 = 1.272 = 1.718
- 0.0592 \times 7.55$).
\textbf{18.} \ce{Cl-} por debajo de las rectas; \ce{Cl2}, en un pequeño
triángulo para $\mathrm{pH} < 1.65$ por encima de 1.446~V; \ce{HClO}, por
encima de la recta de pendiente $-0.0296$ hasta pH 7.55, y \ce{ClO-}, más
allá. La recta \ce{O2}/\ce{H2O}, de 1.23 a 0.40~V, está por debajo de
todas estas fronteras (el diagrama calculado por el programa del capítulo
tiene la misma forma).
\textbf{19.} No: su dominio está por encima de la recta del oxígeno, de
modo que puede oxidar el agua, pero la reacción es lenta. Una piscina
pierde su cloro sobre todo porque este oxida lo que aportan los bañistas
y el aire, y por la luz solar; hay que reponerlo.
\textbf{20.} \ce{Cl2} no tiene dominio allí: se dismuta,
\ce{Cl2 + H2O -> HClO + Cl- + H+}; en medio básico,
\ce{Cl2 + 2OH- -> ClO- + Cl- + H2O}.
\textbf{21.} Haciendo pasar cloro por una disolución de hidróxido de
sodio (la reacción de la pregunta 20).
\textbf{22.} \ce{HClO + Cl- + H+ -> Cl2 + H2O}, una \omterm{def:b1:e-ph-diagrams:disproportionation}{comproporción}.
\textbf{23.} Por debajo de pH 1.6, el hipoclorito y el cloruro
comproporcionan: se forma cloro, que, superada su \omterm{def:b1:precipitation:solubility}{solubilidad}, se
desprende como gas tóxico.
\textbf{24.} Con $c = \qty{0.010}{mol/L}$, el pH 4 está fuera del dominio
de \ce{Cl2}. Pero la frontera depende de $c$: repetir las preguntas
12--14 para una $c$ general da $\mathrm{pH} = 3.34 + \log(2c)$, que llega
a 3.6 para $c = \qty{1}{mol/L}$, el orden de concentración de la lejía.
Mezclar productos concentrados cerca de pH 4 tampoco es seguro.
\textbf{25.} \textbf{pH 1.6}: por encima de él, con
$c = \qty{0.010}{mol/L}$, el cloro disuelto se dismuta en ácido
hipocloroso y cloruro.
\end{solution}
